\chapter{Bevezetés}
\label{ch:intro}

Az informatika területén egyre nehezebb az álláskeresés, mivel egy állásra rengetegen jelentkeznek, ezért a munkáltatók arra vannak kényszerítve, hogy az önéletrajzok alapján kiszűrjék a kevésbé rátermett jelentkezőket. Ezt elsősorban a beküldött önéletrajzok alapján teszik, mivel nem tudnak mindenkit behívni képességeik felmérésére. Ezért kulcsfontosságú egy olyan önéletrajzot leadni, amivel azonnal ki lehet tűnni a többi jelentkező közűl, ezzel elősegítve azt, hogy az interjún is bizonyítani tudjuk tudásunkat, és végül felvételt nyerjünk.

Számtalan mesterséges intelligencia segítségével működő önéletrajz készítő asszisztens létezik, viszont nehéz olyan specializálódott modellt találni, ami a kiberbiztonsági szakterületre jelentkezők igényeit elégítené ki. A szakdolgozatomban bemutatott alkalmazásom azt a célt szolgálja, hogy személyre szabott, szaknyelvet használó önéletrajzot csinál, ami úgy emeli ki a jelöltek releváns ismereteit és kompetenciáit, hogy nemcsak a munkaadó, hanem az ECSF (European Cyber Security Framework) szempontjait is figyelembe veszi.
